\section{ML-KEM-1024 Binding}
\label{sec:mlkem}

\pqcrypt{} invokes ML-KEM-1024 unchanged from FIPS 203 \citep{fips203}: \code{KeyGen}, \code{Encaps(pk)} producing $(K, c)$, and \code{Decaps(sk, c)} producing $K$ (or an implicit-reject pseudorandom value on malformed $c$, per the Fujisaki--Okamoto wrapper). We do not modify the FO transform, do not add a combiner, and do not wrap ML-KEM in any further construction. The shared secret $K$ from \code{Decaps} is fed directly to \code{HKDF-Extract} as IKM (\cref{sec:kdf-nonce}).

\subsection{Tight QROM proof}

ML-KEM-1024's IND-CCA2 security has a tight QROM proof under Module-LWE \citep{hofheinz2017fo,don2022qrom}. \pqcrypt{} inherits this directly. We make no assumption beyond what FIPS 203 already requires.

\subsection{Implementation sources (v1.x, as shipped)}

\begin{center}
\small
\begin{tabular}{@{}lll@{}}
\toprule
\textbf{Language} & \textbf{ML-KEM-1024 source (pinned)} & \textbf{Constant-time claim} \\
\midrule
Python      & \texttt{kyber-py >= 1.0}                                & vendor-documented; NCC Group audit Q2 2024 \\
C           & \textbf{OpenSSL 3.5+ \texttt{EVP\_KEM}} ML-KEM-1024      & OpenSSL constant-time path; FIPS 140-3 candidate \\
Go          & \textbf{Go 1.26 stdlib \texttt{crypto/mlkem}}           & stdlib constant-time \\
Rust        & \texttt{ml-kem = "0.3"} (RustCrypto, \texttt{MlKem1024}) & const-time documented; \texttt{\#![forbid(unsafe\_code)]} \\
C++         & wraps the C library                                     & inherits OpenSSL \\
TypeScript  & \texttt{@noble/post-quantum}                            & const-time documented; pure JS \\
Java        & \textbf{SunJCE} JEP 452 (\emph{not} BouncyCastle)        & JDK stdlib constant-time \\
CPython ext & wraps the C reference (OpenSSL 3.5+ \texttt{EVP\_KEM})    & inherits OpenSSL \\
\bottomrule
\end{tabular}
\end{center}

\textbf{Decisions worth flagging:}
\begin{itemize}
\item \emph{C uses OpenSSL 3.5+ \texttt{EVP\_KEM}, not \texttt{mlkem-native}.} The earlier draft of this whitepaper named \texttt{mlkem-native} \citep{cryspen2025mlkem}; the PQCK-00005 release note records the decision to switch. OpenSSL has a smaller audit surface for our deployment (we already link \texttt{libcrypto} for the AEADs), a FIPS 140-3 validation path, and ubiquitous packaging. \texttt{mlkem-native} is \emph{not} vendored, linked, or retained by this project in any capacity; an embedded target that cannot link OpenSSL would have to bring its own ML-KEM source, which requires the security review mandated by the supply-chain rules in \texttt{tessera-seal/CLAUDE.md} and is outside v1.x scope.
\item \emph{Java uses SunJCE (JEP 452), not BouncyCastle.} The JEP 452 \texttt{KEM} API landed in JDK 24 and provides ML-KEM-1024 as a first-class \texttt{KeyPairGenerator} algorithm. BouncyCastle 1.79+ remains a viable alternative but is not the v1.1 reference; this avoids the CVE-2025-1432 surface (\cref{sec:side-channels}) in the default vendoring story.
\item \emph{Go uses stdlib \texttt{crypto/mlkem} (Go 1.26+), not \texttt{filippo.io/mlkem}.} The earlier external library was the v1 alpha intent; by the time v1.1 shipped, ML-KEM had landed in the Go stdlib and the dependency was eliminated.
\end{itemize}

The v1 alpha language set (Python, C, Go, Rust) is the minimum that exercises the spec across a readable reference (Python), a canonical ABI (C), a stdlib-heavy ecosystem (Go), and the strongest constant-time toolchain (Rust). C++, TypeScript, Java, and CPython ext joined in v1.1; see Appendix F for the release timeline.

\subsection{The ``Kyber'' naming ban}

ML-KEM (FIPS 203, August 2024) is wire-incompatible with the Kyber-r3 specification it descends from. \pqcrypt{}'s public APIs MUST NOT expose the token ``Kyber'' even where the upstream library uses it. Rename at the binding layer (e.g., \texttt{kyber-py.ML\_KEM\_1024} $\to$ \texttt{tessera\_seal.MLKem1024}). The same discipline applies to the hybrid suite's classical half: only ``X25519'' appears in user-facing identifiers, never ``Curve25519'' or ``ECDH.''
